\exercisesection{1}

\begin{enumerate}
\def\labelenumi{\arabic{enumi}.}
\item
  设 \(d\) 为一个在 \(\mathbf{Z}\) 中不被平方数整除的有理整数。证明：域 \(\mathbf{K} = \mathbf{Q}(\sqrt{d})\) 中在 \(\mathbf{Z}\) 上整的元素形如 \(a + b\sqrt{d}\)（其中 \(a \in \mathbf{Z}\)、\(b \in \mathbf{Z}\)），若 \(d - 1 \not\equiv 0 \pmod{4}\)；而形如 \((a + b\sqrt{d})/2\)（其中 \(a \in \mathbf{Z}\)、\(b \in \mathbf{Z}\)，且 \(a\) 与 \(b\) 同为偶数或 \(a\) 与 \(b\) 同为奇数），若 \(d - 1 \equiv 0 \pmod{4}\)（参见《代数》第七章 § 1 习题 8）。由此推出 \(\mathbf{A} = \mathbf{Z}[\sqrt{5}]\) 不是整闭的，并给出 \(\mathbf{Q}[\sqrt{5}]\) 中一个在 \(\mathbf{A}\) 上整的元素的例子，但它在 \(\mathbf{A}\) 的分式域上的极小多项式的系数不属于 \(\mathbf{A}\)。
\item
  设 K 为交换域，A 为多项式代数 \(K[X, Y]\) 中由单项式 \(Y^{k}X^{k+1} (k \geqslant 0)\) 生成的子 K-代数。证明：XY 使得 A{[}XY{]} 包含于一个有限生成的 A-模，但 XY 不在 A 上整。
\item
  给出交换域 K 的扩张的一个无穷序列 \((\mathbf{K}_{n})\) 的例子，其中每个扩张在 K 上有限次，使得 K-代数 \(\prod_{n} K_{n}\) 不在 K 上代数。
\item
  在矩阵环 \(\mathbf{M}_{2}(\mathbf{Q})\) 中给出两个在 Z 上整的元素的例子，但它们的和与积都不在 Z 上整（考虑形如 \(I + N\) 的矩阵，其中 N 幂零）。
\item
  设 A 为交换环，B 为包含 A 的交换环，x 为 B 的一个可逆元。证明 \(A[x] \cap A[x^{-1}]\) 的每个元素都在 A 上整。（若 \(y = a_{0} + \cdots + a_{p}x^{-p} = b_{0} + \cdots + b_{q}x^{q}\)，其中 \(a_{i}\) 与 \(b_{j}\) 属于 A，证明由 1，\(x, \ldots, x^{p+q+1}\) 生成的 B 的子 A-模是一个忠实 A{[}y{]}-模。）
\item
  设 A 为交换环，B 为交换 A-代数；假设 B 的极小素理想个数有限；设 \(p_{i}\)（\(1 \leqslant i \leqslant n\)）为这些理想（注意若 B 是 Noether 环则此假设成立）。元素 \(x \in B\) 在 A 上整，当且仅当它的每个典范像 \(x_{i}\)（在 \(B/p_{i}\) 中）都在 A 上整。（若此条件成立，先证明 x 在 B 的由幂零根 \(N = \bigcap_{i} p_{i}\) 生成的子代数上整，并利用 N 的元素皆幂零这一事实。）
\item
  给出一个约化环 A 的例子，它在其全分式环中整闭，但不是整环之积（取域 K 在无穷积 \(\prod_{n} K_{n}\)（即 K 的适当代数扩张之积）中的整闭包）。
\item
  设 A 为交换环，B 为有限生成 A-代数。证明存在一个作为自由 A-模的有限 A-代数 C 以及一个满射 A-同态 \(C \rightarrow B\)。
\item
  证明商整环 \(\mathbf{Z}[\mathbf{X}] / (\mathbf{X}^2 +4)\) 不是整闭的，尽管 \(\mathbf{Z}[\mathbf{X}]\) 是整闭的。
\item
  设 A 为完全整闭整环。证明：对每个集合 I，多项式整环 \(A[X_{i}]_{i \in I}\) 与形式幂级数整环 \(A[[X_{i}]]_{i \in I}\)（《代数》第四章 §5 习题 1）都是完全整闭的。（利用第 4 节命题 4 及下面的引理：对 I 的每个有限子集 J 及所有 \(P \in A[[X_{i}]]_{i \in J}\)，设 \(P_{J}\) 为 \(A[[X_{i}]]_{i \in J}\) 中由 P 出发、把 P 中所有 \(X_{i}\)（指 \(i \notin J\) 者）换为 0 而得的形式幂级数；若 P、Q 是两个形式幂级数，使得 \(P_{J}\) 被 \(Q_{J}\) 整除（在 \(A[[X_{i}]]_{i \in J}\) 中，对所有 J），则 P 在 \(A[[X_{i}]]_{i \in I}\) 中被 Q 整除。为此注意：若 J、\(J'\) 是 I 的两个有限子集，满足 \(J \subset J'\)，且 \(P_{J'} = L'Q_{J'}\) 与 \(P_{J} = LQ_{J}\)，则 L 由 \(L'\) 出发、把 \(L'\) 中所有 \(X_{i}\)（指 \(i \in J' - J\) 者）换为 0 而得到。）
\item
  \begin{enumerate}
  \def\labelenumii{(\alph{enumii})}
  \tightlist
  \item
    设 R 为整环，\((\mathbf{A}_{\alpha})\) 为 R 的一个右定向的子环族，A 为诸 \(A_{\alpha}\) 的并。证明：若每个 \(A_{\alpha}\) 都整闭，则 A 也整闭。
  \end{enumerate}
\end{enumerate}

\begin{enumerate}
\def\labelenumi{(\alph{enumi})}
\setcounter{enumi}{1}
\tightlist
\item
  设 K 为交换域，\(\mathbf{R} = \mathbf{K}(\mathbf{X}, \mathbf{Y})\) 为 K 上两个未定元的有理函数域。对每个整数 n \textgreater{} 0，设 \(A_{n}\) 表示 R 的子环 \(\mathbf{K}[\mathbf{X}, \mathbf{Y}, \mathbf{Y}\mathbf{X}^{-n}]\)。证明诸 \(A_{n}\) 都是完全整闭的，但它们的并 A 不是。（为证形如 \(\mathbf{P}(\mathbf{X}, \mathbf{Y}) \cdot \mathbf{X}^{-ns}\) 的元素（其中 \(P \in K[X, Y]\)）属于 \(A_{n}\)，证明其必要且充分条件是 P 属于 \(S^{s}\)，其中 S 表示 \(K[X, Y]\) 中由 \(X^{n}\) 与 Y 生成的理想。）
\end{enumerate}

¶ * 12. 设 A 为复变量的（取复值的）整函数环。它是一个整环，其分式域是复变量的亚纯函数域。

\begin{enumerate}
\def\labelenumi{(\alph{enumi})}
\item
  证明 \(\mathbf{A}\) 是完全整闭的。
\item
  对所有 \(x \in \mathbf{A}\)，设 \(Z(x)\) 为 \(x\) 的零点集（它是 \(\mathbf{C}\) 的闭子集，当 \(x \neq 0\) 时是离散的（从而可数））。若 \(\mathfrak{a}\) 是 \(\mathbf{A}\) 的一个异于 \(\mathbf{A}\) 与 (0) 的理想，证明诸 \(Z(x)\)（对 \(x \in \mathfrak{a}\)）构成一个滤子基 \(\mathfrak{F}_{\mathfrak{a}}\)。（若 \(x\) 与 \(y\) 是两个没有公共零点的整函数，试利用 Mittag-Leffler 定理证明可以写出 \(1 / xy = u / x + v / y\)，其中 \(u\) 与 \(v\) 属于 \(\mathbf{A}\)，并由此推出 \(Ax + Ay = A\)。）反之，对每个滤子 \(\mathfrak{F}\)（在 \(\mathbf{C}\) 上、且 \(\mathbf{C}\) 的某个闭离散子集属于它），集合 \(\mathfrak{a}(\mathfrak{F})\)（由 \(x \in \mathbf{A}\) 中使 \(Z(x) \in \mathfrak{F}\) 者组成）是 \(\mathbf{A}\) 的一个理想。证明：\(\mathfrak{F}_{\mathfrak{a}(\mathfrak{F})} = \mathfrak{F}\) 对每个滤子 \(\mathfrak{F}\)（\(\mathbf{C}\) 的某个闭离散子集属于它）成立，且有 \(\mathfrak{a}(\mathfrak{F}) \supset \mathfrak{a}\) 对每个理想 \(\mathfrak{a}\)（\(\mathbf{A}\) 的、异于 \(\mathbf{A}\) 与 (0) 的）成立。
\item
  若 \(\mathfrak{p} \neq (0)\) 是 \(\mathbf{A}\) 的素理想，证明 \(\mathfrak{F}_{\mathfrak{p}}\) 是超滤子。（限于 \(\mathfrak{F}_{\mathfrak{p}}\) 中所有集合都是无穷集的情形，证明：若 \(\mathbf{M} \in \mathfrak{F}_{\mathfrak{p}}\) 是两个无公共点的无穷集 \(\mathbf{M}'\)、\(\mathbf{M}''\) 之并，则其中之一属于 \(\mathfrak{F}_{\mathfrak{p}}\)。）反之，对每个超滤子 \(\mathfrak{U}\)（在 \(\mathbf{C}\) 上、且 \(\mathbf{C}\) 的某个闭离散子集属于它），\(\mathfrak{a}(\mathfrak{U})\) 是 \(\mathbf{A}\) 的一个极大理想；试得出其逆。
\item
  设 \(\mathfrak{U}\) 为 \(\mathbf{C}\) 上的一个超滤子，它是在典范单射 \(\mathbf{Z} \to \mathbf{C}\) 下、一个 \(\mathbf{Z}\) 上的非平凡超滤子的像。证明存在非极大的素理想 \(\mathfrak{p}\)（\(\mathbf{A}\) 的）使 \(\mathfrak{F}_{\mathfrak{p}} = \mathfrak{U}\)。（对所有 \(x \in \mathfrak{a}(\mathfrak{U})\) 及所有 \(n \in \mathbf{Z}\)，设 \(\omega_x(n)\) 为 \(x\) 在点 \(n\) 处的阶（\(\omega_x(n) = 0\)，若 \(x(n) \neq 0\)）；考虑集合 \(\mathfrak{p}\)：它由诸 \(x \in \mathfrak{a}(\mathfrak{U})\)（使 \(\lim_{\mathfrak{u}} \omega_x = +\infty\) 者）组成。）
\item
  用 (d) 的记号，设 m 为极大理想 \(\mathfrak{a}(\mathfrak{U})\)。证明整环 \(A_{m}\) 不是完全整闭的（考虑亚纯函数 \(1/(\sin\pi z)\)）。
\end{enumerate}

¶ 13. 设 A 为局部整环，K 为其分式域。考虑下面两个性质：

\begin{enumerate}
\def\labelenumi{(\roman{enumi})}
\item
  A 是 Hensel 的（第三章 § 4 习题 3）。
\item
  对每个有限扩张 \(\mathbf{L}\)（\(\mathbf{K}\) 的），整闭包 \(\mathbf{A}'\)（即 \(\mathbf{A}\) 在 \(\mathbf{L}\) 中的整闭包）是局部环。
\end{enumerate}

证明 (i) 蕴含 (ii)，且当 A 整闭时其逆为真。（为证 (i) 蕴含 (ii)，注意 A' 是一个由有限 A-代数组成的右定向族的并，并利用 Hensel 环的定义。为证当 A 整闭时 (ii) 蕴含 (i)，证明：若 P ∈ A{[}X{]} 是首一不可约多项式，且 f: A → k 是 A 到其剩余域 k 上的典范同态，则 \(\bar{f}(P)\) 在 K{[}X{]} 中不可约，并考虑域 L = K{[}X{]}/(f)。）

若条件 (ii) 满足，则 A 在 K 的任何扩张中的整闭包 A' 都是局部环。考察完备 Hausdorff 局部环的情形；若 A 是完备 Hausdorff Noether 局部环，则每个含于 A' 且作为有限生成 A-模的 A-代数都是完备 Hausdorff 局部环。

\begin{enumerate}
\def\labelenumi{\arabic{enumi}.}
\setcounter{enumi}{13}
\tightlist
\item
  设 A 为完全整闭整环，K 为其分式域。证明对 K 的每个扩张 L，A 在
\end{enumerate}

L 中的整闭包 A' 是完全整闭整环。（化归到 L 是 K 的次数为 m 的有限拟 Galois 扩张的情形：若 \(x \in L\) 使得存在 \(d \in A'\) 使 \(dx^{n} \in A'\)（对所有 \(n \geqslant 0\)），先证明可设 \(d \in A\)，然后推出：若 \(c_{i}\)（\(1 \leqslant i \leqslant m\)）是 x 在 K 上的极小多项式的系数，则 \(d^{m}c_{i}^{n} \in A\)（对所有 \(n \geqslant 0\)）。）

\begin{enumerate}
\def\labelenumi{\arabic{enumi}.}
\setcounter{enumi}{14}
\tightlist
\item
  设 K 为特征 0 的交换域，它包含所有单位根，并且存在 K 的 Galois 扩张，其 Galois 群不可解（例如参见《代数》第五章附录 I 第 1 节命题 2）。设 \(\Omega\) 为 K 的一个代数闭包，并设 \(\Omega_0\) 为满足下列条件的元素 \(z \in \Omega\) 组成的集：由 \(z\) 在 \(\Omega\) 中生成的 K 的 Galois 扩张具有可解的 Galois 群。
\end{enumerate}

\begin{enumerate}
\def\labelenumi{(\alph{enumi})}
\item
  证明 \(\Omega_0\) 是 \(\Omega\) 的一个异于 \(\Omega\) 的子群（参见《代数》第五章 § 10 第 4 节定理 1）。
\item
  设 A 为 \(\Omega[X]\) 的子环，由满足 \(\mathrm{P}(0) \in \Omega_{0}\) 的多项式 P 组成。证明 A 不是整闭的，但在其分式域 \(\Omega(X)\) 中，每个有某次幂属于 A 的元素本身属于 A。
\end{enumerate}

\begin{enumerate}
\def\labelenumi{\arabic{enumi}.}
\setcounter{enumi}{15}
\item
  设 B 为环，A 为 B 的子环，f 为 A 的理想，即 A-模 B/A 的零化子。设 U 为 \(\mathbf{X} = \operatorname{Spec}(\mathbf{A})\) 中闭集 \(\mathbf{V}(\mathfrak{f})\) 的余集，并设 \(u = {}^{a}\rho\)，其中 \(\rho: A \to B\) 是典范单射。证明 u 在 \({}^{-1}(U)\) 上的限制是 \({}^{-1}(U)\) 到 U 上的同胚（对每个元素 \(g \in f\)，考虑环 \(A_{g}\) 的谱，它与开集 \(X_{g} \subset U\) 等同）。
\item
  设 K 为域，B 为多项式环 \(K[X_{n}]_{n\in N}\)，A 为 B 的子环，由单项式 \(X_{n}^{2}\) 与 \(X_{n}^{3}\)（对所有 \(n\in N\)）生成；证明 B 是 A 的整闭包，且 B 在 A 中的导子化为 0；但若 \(S=A-\{0\}\)，则 \(S^{-1}B\) 在 \(S^{-1}A\) 中的导子不化为 0，且 \(S^{-1}A\) 是整闭的。
\item
  设 A 为整环，K 为其分式域，M 为有限生成 A-模，u 为 M 的一个自同态，\(u \otimes 1\) 为有限维 K-向量空间 \(M_{(K)} = M \otimes_{A} K\) 上相应的自同态。证明 \(u \otimes 1\) 的特征多项式的系数在 A 上整（利用第 1 节定理 1 的条件 \((E_{\mathrm{III}})\)）。
\item
  \begin{enumerate}
  \def\labelenumii{(\alph{enumii})}
  \tightlist
  \item
    设 A 为整闭整环，K 为其分式域，L 为 K 的可分代数扩张，x 为 L 中在 A 上整的一个元素，\(K'\) 为域 \(\mathrm{K}(x) \subset \mathrm{L}\)，F 为 x 在 K 上的极小多项式。若 x 在 K 上的次数为 n，证明 A 在 \(K'\) 中的整闭包包含在由元素 \(1/\mathrm{F}'(x)\)、\(x/\mathrm{F}'(x)\)、\(\ldots\)、\(x^{n-1}/\mathrm{F}'(x)\)（都属于 \(K'\)）生成的 A-模之中。（若 \(z \in K'\) 在 A 上整，写下
  \end{enumerate}
\end{enumerate}

\[
z \mathrm{F} ^ {\prime} (x) = g (x), \quad \text { where } \quad g (\mathbf {X}) = b _ {0} + b _ {1} \mathbf {X} + \dots + b _ {n - 1} \mathbf {X} ^ {n - 1} \in \mathbf {K} [ \mathbf {X} ].
\]

借助 Lagrange 插值公式证明

\[
g (\mathbf {X}) = \operatorname{Tr} _ {\mathbf {K} ^ {\prime} [ \mathbf {X} ] / \mathbf {K} [ \mathbf {X} ]} \left(z. \frac {\mathrm{F} (\mathbf {X})}{\mathbf {X} - x}\right) \cdot)
\]

\begin{enumerate}
\def\labelenumi{(\alph{enumi})}
\setcounter{enumi}{1}
\tightlist
\item
  设 A 为整闭整环，K 为其分式域，p 为 K 的特征指数。假设子环 \(A^{1/p}\)（\(K^{1/p}\) 中由 A 的元素的 p 次根组成的子环）是有限生成 A-模。证明对 K 的每个有限扩张 E，A 在 E 中的整闭包是有限生成 A-模（化归到 E 是拟 Galois 扩张的情形）。
\end{enumerate}

¶ * 20. 设 \(\mathbf{K}_0\) 为特征 \(p > 0\) 的完满域，\(\mathbf{K}\) 为有理函数域 \(\mathrm{K}_0(\mathrm{X}_n)_{n \in \mathbb{N}}\)，\(\mathbf{B}\) 为形式幂级数环 \(\mathrm{K}[[\mathrm{Z}]]\)，其中 \(\mathbf{Z}\) 是一个未定元；若 \(m\) 是 \(B\) 的极大理想，则 \(B\) 对 \(m\) -进拓扑是 Hausdorff 且完备的，并且是一个离散赋值环。设 \(L = K((Z))\) 为 \(B\) 的分式域，\(E_0\) 为 \(L\) 的子域，由 \(L^p\) 与元素 \(Z\) 和 \(X_n\)（对 \(n \in N\)）生成。

\begin{enumerate}
\def\labelenumi{(\alph{enumi})}
\item
  证明元素 \(c = \sum_{n=0}^{\infty} \mathbf{X}_n \mathbf{Z}^n\)（属于 \(L\)）不属于 \(E_0\)。
\item
  设 E 为 L 的包含 \(E_{0}\) 而不包含 c 的子域所成之集的一个极大元；记 \(A = B \cap E\)。证明 A 是一个离散赋值环（从而是 Noether 的），其分式域为 E；若 \(m'\) 是 A 的极大理想，则 \(m'^{k} = m^{k} \cap A\)，且 A 对 m-进拓扑在 B 中稠密；\(L = K(c)\) 是 K 的一个次数为 p 的扩张，B 是 A 在 L 中的整闭包，但 B 不是有限生成 A-模（利用第三章 §3 习题 9 (b)）。*
\end{enumerate}

\begin{enumerate}
\def\labelenumi{\arabic{enumi}.}
\setcounter{enumi}{20}
\tightlist
\item
  \begin{enumerate}
  \def\labelenumii{(\alph{enumii})}
  \tightlist
  \item
    设 \(K_{0}\) 为特征 p \textgreater{} 0 的完满域，L 为有理函数域 \(\mathrm{K}_{0}(\mathrm{X}_{n})_{n\in\mathbf{N}}\)，\(A'\) 为形式幂级数环 \(L[[Y,Z]]\)（其中 Y 与 Z 是两个未定元），A 为张量积环 \(K[[Y,Z]]\otimes_{K}L\)，这里记 \(K = L^{p} \subset L\)；A 是一个 Noether 局部环，它不完备，其完备化与 \(A'\) 等同（第三章 §3 习题 17），且 \(A'\) 在 A 上整。设 \(c_{n}\) 表示元素
  \end{enumerate}
\end{enumerate}

\[
\mathrm{X} _ {n} \mathrm{Y} + \mathrm{X} _ {n + 1} \mathrm{YZ} + \mathrm{X} _ {n + 2} \mathrm{YZ} ^ {2} + \dots
\]

（属于 A'）。设 B 为 A' 的由 A 与诸元素 \(c_{n}\)（\(n \geqslant 0\)）生成的子环。若 C = A{[}c₀{]}，则 C 是 Noether 的，且环 B（它不整闭）含于 C 的整闭包并包含 C；但证明 B 不是 Noether 的（注意 \(c_{n}\) 不属于 B 的由诸 \(c_{i}\)（其下标 i \textless{} n）生成的理想）。

\begin{enumerate}
\def\labelenumi{(\alph{enumi})}
\setcounter{enumi}{1}
\tightlist
\item
  设 \(p = 2\)；\(K_0\)、\(K\) 与 \(L\) 的意义与 (a) 中相同，这次考虑三个未定元的形式幂级数环 \(A' = L[[Y, Z, T]]\)，以及子环
\end{enumerate}

\[
\mathrm{A} = \mathrm{K} [ [ \mathrm{Y}, \mathrm{Z}, \mathrm{T} ] ] \otimes_ {\mathrm{K}} \mathrm{L},
\]

并写为

\[
b = \mathrm{Y} \sum_ {n \geqslant 0} \mathrm{X} _ {2 n} \mathrm{T} ^ {n} + \mathrm{Z} \sum_ {n \geqslant 0} \mathrm{X} _ {2 n + 1} \mathrm{T} ^ {n}.
\]

证明环 A{[}b{]} 是 Noether 的，但它的整闭包 B 不是 Noether 环。（由于 \(b^{2} \in A\)，A{[}b{]} 的分式域的每个元素形如 \(P + Qb\)，其中 P 与 Q 是诸 \(X_{k}\) 中有限个的线性组合，系数在 K{[}{[}Y, Z, T{]}{]} 的分式域中；为使这样一个元素在 A{[}b{]} 上整，证明其必要且充分条件是它属于 A'。证明 B 由诸元素

\[
b _ {i} = \mathrm{Y} \sum_ {n \geqslant 0} \mathrm{X} _ {2 (i + n)} \mathrm{T} ^ {n} + \mathrm{Z} \sum_ {n \geqslant 0} \mathrm{X} _ {2 (i + n) + 1} \mathrm{T} ^ {n} \quad \text { for } \quad i \geqslant 0
\]

生成，并像 (a) 那样完成论证。）

\begin{enumerate}
\def\labelenumi{\arabic{enumi}.}
\setcounter{enumi}{21}
\item
  设 A 为整闭整环，K 为其分式域；证明对 K 的每个扩张 L，存在 L 的一个整闭子整环 B，其分式域是 L 且满足 \(B \cap K = A\)（把 L 看作 K 的一个纯超越扩张 \(E = K(X_t)_{t \in I}\) 的代数扩张，从而利用习题 11 (a) 化归到 L 是 K 的代数扩张的情形）。
\item
  \begin{enumerate}
  \def\labelenumii{(\alph{enumii})}
  \tightlist
  \item
    设 K 为交换域，n 为与 K 的特征互素的整数，P 为 K{[}X{]} 中一个只有单根（在 K 的一个代数闭包中）的多项式。若记 \(f(\mathbf{X}, \mathbf{Y}) = \mathbf{Y}^{n} - \mathbf{P}(\mathbf{X})\)，证明：当 K 包含 n 次单位根时，整环 \(\mathbf{A} = \mathbf{K}[\mathbf{X}, \mathbf{Y}] / (f)\) 是整闭的，并且在其分式域中含有 K(X) 的由 f 的一个根 y 生成的可分扩张 L（f 视为 K(X){[}Y{]} 中的多项式）。（证明：若 L 的元素 z 在 K{[}X{]} 上整，则它属于 A；为此，把 z 写成形式
  \end{enumerate}
\end{enumerate}

\[
z = \sum_ {i = 0} ^ {n - 1} a _ {i} (\mathbf {X}) y ^ {i},
\]

其中 \(a_{i}(\mathbf{X}) \in \mathbf{K}(\mathbf{X})\)。证明诸元素 \(a_{i}(\mathbf{X})y^{i} (0 \leqslant i \leqslant n - 1)\) 在 \(\mathbf{K}[\mathbf{X}]\) 上整，并推出诸 \((a_{i}(\mathbf{X}))^{n}(\mathbf{P}(\mathbf{X}))^{i}\) 是 \(\mathbf{K}[\mathbf{X}]\) 的元素；从而诸 \(a_{i}(\mathbf{X})\) 也是。

\begin{enumerate}
\def\labelenumi{(\alph{enumi})}
\setcounter{enumi}{1}
\tightlist
\item
  设 K 为特征 \(p \neq 2\) 的非完满域；设 \(a \in K\) 为一个不属于 \(K^{p}\) 的元素，设 E 为域 \(\mathrm{K}(a^{1/p})\)，并记 \(\mathrm{P}(\mathrm{X}) = \mathrm{X}^{p} - a, f(\mathrm{X}, \mathrm{Y}) = \mathrm{Y}^{2} - \mathrm{P}(\mathrm{X})\)。若 A 为整闭整环 \(\mathrm{K}[\mathrm{X}, \mathrm{Y}]/(f)\)，证明 \(B = E \otimes_{K} A\) 是整环但不整闭（考虑 B 的分式域中的元素 \(y/(\mathrm{X} - a^{1/p})\)）。
\end{enumerate}

\begin{enumerate}
\def\labelenumi{\arabic{enumi}.}
\setcounter{enumi}{23}
\tightlist
\item
  设 \(\Delta\) 为一个以加法记的无扭交换群。设 A 为交换环，B 为交换 A-代数，\(h\colon\mathbf{A}\to\mathbf{B}\) 为定义 B 上代数结构的环同态。那么 \(h^{(\Delta)}\colon\mathbf{A}^{(\Delta)}\to\mathbf{B}^{(\Delta)}\)（《代数》
\end{enumerate}

第二章 § 11 习题 1）是代数 \(\mathbf{A}^{(\Delta)}\)（群 \(\Delta\) 关于 \(\mathbf{A}\) 的代数）到代数 \(\mathbf{B}^{(\Delta)}\)（群 \(\Delta\) 关于 \(\mathbf{B}\) 的代数）的一个（环）同态。设 \(\mathbf{P} = \sum_{\alpha \in \Delta} b_{\alpha} \otimes \mathbf{X}^{\alpha}\) 为 \(\mathbf{B}^{(\Delta)}\) 的一个元素。证明：\(\mathbf{P}\) 在 \(\mathbf{A}^{(\Delta)}\) 上整，当且仅当每个元素 \(b_{\alpha} \in \mathbf{B}\) 都在 \(\mathbf{A}\) 上整。（化归到 \(\Delta\) 是有限生成的情形，再化归到 \(\Delta = \mathbf{Z}\) 的情形。）证明：若 \(\mathbf{A}\) 是整环且整闭，则 \(\mathbf{A}^{(\Delta)}\) 也整闭。

\begin{enumerate}
\def\labelenumi{\arabic{enumi}.}
\setcounter{enumi}{24}
\tightlist
\item
  设 \(\Delta\) 为全序交换群。把第 8 节的命题 20 与 21 推广到 \(\Delta\) 型分次环（对命题 20 利用习题 24；对命题 21，把第 8 节的引理 4 推广到 \(\Delta\) 有限生成的情形，然后取正向极限证明命题 21）。
\end{enumerate}

¶ 26. (a) 设 A 为满足下列条件的整环：对 A 中每个 \(a \neq 0\)，理想 \(Aa\) 是一个族 \((q_t)\)（由 \(Aa\) 关于素理想 \(p_t\)（在包含 \(a\) 的素理想中极小者）的饱和组成）之交，且诸整环 \(A_{p_t}\) 是整闭的（相应地，完全整闭的）。证明 A 是整闭的（相应地，完全整闭的）。

\begin{enumerate}
\def\labelenumi{(\alph{enumi})}
\setcounter{enumi}{1}
\item
  设 A 为强 Lasker 局部环（第四章 § 2 习题 28），m 为其极大理想。假设对某个 \(a \in \mathbf{A}\)，m 是与 A/Aa 弱伴随的一个浸入素理想（第四章 § 1 习题 17）。证明 Aa: m 含于每个对与 A/Aa 弱伴随的素理想 p ≠ m 都准素的理想之中（对这样一个准素理想 q，考虑传递子 q: m，参见第四章 § 2 习题 30）。若进一步设 A 是整环，证明关系式 p.~(Aa: m) = Aa（对在包含 Aa 的素理想中极小的素理想 p）不可能成立（若 q 是 A 关于 p 的饱和，注意上述关系将蕴含在 A \(_{p}\) 中 qA \(_{p}\) = qpA \(_{p}\)，并借助第四章 § 2 习题 29 (d) 完成证明）。
\item
  证明：若 A 是强 Lasker 整环，且存在 A 的元素 \(a \neq 0\) 使得有与 A/Aa 弱伴随并且浸入的素理想，则 A 不是完全整闭的。（化归到 (b) 中考虑的情形；用 (b) 的记号，那时存在 \(b \in \mathbf{A}a\) : m 使得 \(b \notin \mathbf{A}a\) 且 \(bp \subset am\)；对 \(n\) 用归纳法推出 \(b^n p \subset a^n m\)（对所有 \(n\)）。）
\item
  特别地，由 (a) 与 (c) 推出：Noether 整环是整闭的，当且仅当对 A 中所有 \(a \neq 0\)，与 A/Aa 伴随的素理想 \(\mathfrak{p}\) 都不是浸入的，且 \(\mathrm{A}_{\mathfrak{p}}\) 是整闭的（参见第七章 § 1 第 4 节命题 8）。
\end{enumerate}

\begin{enumerate}
\def\labelenumi{\arabic{enumi}.}
\setcounter{enumi}{26}
\tightlist
\item
  设 A 为整闭但非完全整闭的整环，K 为其分式域；设 \(a \neq 0\) 为 A 的一个不可逆元素，对它存在 A 中的 \(d \neq 0\) 使 \(da^{-n} \in A\)（对每个整数 \(n \geqslant 0\)）。
\end{enumerate}

证明在形式幂级数域 \(\mathbf{K}((\mathbf{X}))\) 中存在一个形式幂级数 \(\sum_{k=1}^{\infty} u_k \mathbf{X}^k = f(\mathbf{X})\)，满足方程

\[
(f (\mathbf {X})) ^ {2} - a f (\mathbf {X}) + \mathbf {X} = 0,
\]

从而它在 A{[}{[}X{]}{]} 上整，属于 A{[}{[}X{]}{]} 的分式域但不属于 A{[}{[}X{]}{]}，故 A{[}{[}X{]}{]} 不是整闭的。

\begin{enumerate}
\def\labelenumi{\arabic{enumi}.}
\setcounter{enumi}{27}
\item
  设 \(k\) 为域，\(\mathbf{A}_0\) 为两个无穷未定元族的多项式环 \(k[\mathrm{X}_n, \mathrm{Y}_n]_{n \in \mathbf{Z}}\)。设 \(G\) 为一个无穷单生成群（从而同构于 \(\mathbf{Z}\)），并设 \(\sigma\) 为 \(\mathcal{G}\) 的一个生成元。\(G\) 定义为在 \(\mathbf{A}_0\) 上作用，其条件为 \(\sigma(\mathrm{Y}_n) = \mathrm{Y}_n\) 与 \(\sigma(\mathrm{X}_n) = \mathrm{X}_{n+1}\)（对所有 \(n \in \mathbf{Z}\)）。设 \(\Im\) 为 \(\mathbf{A}_0\) 中由元素 \(\mathrm{Y}_n(\mathrm{X}_n - \mathrm{X}_{n+1})\)（对所有 \(n \in \mathbf{Z}\)）生成的理想，并设 \(A\) 为商环 \(A_0 / \Im\)；由于 \(\sigma(\Im) \subset \Im\)，群 \(\mathcal{G}\) 通过过渡到商也在 \(A\) 上作用。设 \(S\) 为 \(A\) 中由诸 \(Y_n\) 在 \(A\) 中的典范像生成的乘性子集，于是 \(S\) 由在 \(G\) 下不变的元素组成。证明 \(S^{-1}(A^{\mathcal{G}}) \neq (S^{-1}A)^{\mathcal{G}}\)。
\item
  设 k 为特征 \(\neq2\) 的域，A 为一个未定元的多项式环 k{[}T{]}，a、b 为 k 中两个元素，满足 \(a\neq0\)、\(b\neq0\) 且 \(a\neq b\)。设 \(\mathbf{K}=k(\mathbf{T})\) 为 A 的分式域，L、M 为 K 的扩张，分别由 \(\mathbf{X}^{2}-\mathbf{T}(\mathbf{T}-a)\) 的一个根和 \(\mathbf{X}^{2}-\mathbf{T}(\mathbf{T}-b)\) 的一个根添加到 K 而得。设 B、C 分别为 A 在 L 与 M 中的整闭包，它们是有限生成自由 A-模。证明 \(L\otimes_{K}M\) 是一个域，是 K 的可分有限扩张，但 \(B\otimes_{A}C\)（它与 \(L\otimes_{K}M\) 的一个子环典范地等同）不是 A 在 \(L\otimes_{K}M\) 中的整闭包。
\end{enumerate}
% semantic-fragments: legacy-input-seam
\relax{}
